\documentclass[10pt,a4paper,twocolumn]{article}
\usepackage{kotex}
\usepackage[margin=18mm,columnsep=7mm]{geometry}
\usepackage{booktabs}
\usepackage{multirow}
\usepackage{microtype}
\usepackage{graphicx}
\usepackage{xcolor}
\usepackage{url}
\usepackage[authoryear,round]{natbib}
\usepackage{hyperref}
\hypersetup{colorlinks=true,linkcolor=black,citecolor=black,urlcolor=blue}
\setlength{\parindent}{1em}
\setlength{\parskip}{0pt}
\setlength{\emergencystretch}{3em}
\Urlmuskip=0mu plus 1mu
\sloppy
\raggedbottom

\newcommand{\dataset}{K-FNSPID v3}
\newcommand{\system}{Hana Montana AI(KF-DeBERTa + K-FNSPID)}
\newcommand{\model}{KF-DeBERTa}
\renewcommand{\abstractname}{초록}
\renewcommand{\tablename}{표}
\renewcommand{\refname}{참고문헌}

\title{K-FNSPID v3: 미래 정보 누수를 통제한\\한국 금융 뉴스·공시·시장반응 데이터셋}
\author{최성현\\
한국공학대학교 SW대학 컴퓨터공학부 소프트웨어전공\\
4학년 학부생}
\date{2026년 7월}

\begin{document}
\maketitle

\begin{abstract}
본 연구는 한국 상장시장의 뉴스 524,696건과 전자공시 25,966건을 파일 기반 일별 시세 10,691,998행에 연결한 금융 자연어처리 자원 \dataset{}를 제안한다. 동결 배포본은 문서 550,662건, 문서-종목 관계 819,772건, 시장반응 표본 398,942건과 사건 혼입 제거 후 대표 표본 130,566건으로 구성된다. 발표 시각을 한국 거래 세션에 맞춰 정규화하고, 유사 사건을 군집화하며, 동일 종목·거래일의 서로 다른 사건 충돌을 제외하고, 7일 embargo를 둔 시간순 분할로 평가한다. 감성분류에서 \model{} LoRA는 균형 공개 시험셋 macro-F1을 KR-FinBERT-SC의 0.7272에서 0.8850으로 높인다. 순서형 시장반응에서는 Train 전용 TF-IDF 기준선 대비 macro-F1을 0.3429에서 0.3820으로, quadratic weighted kappa를 0.3141에서 0.4694로 높였으며 거래일 군집 bootstrap 신뢰구간도 0을 넘는다. 그러나 개선은 뉴스에 집중되고 공시 하위집합의 macro-F1은 하락한다. 공시 의미 중요도는 별도 910건 코드북 평가셋에서 검증한다. 본 데이터와 모델은 의미 중요도, 사후 시장반응, 인과효과를 구분하며 투자수익을 주장하지 않는다.
\end{abstract}

\section{서론}

금융 모니터링 시스템은 문장의 긍정·부정뿐 아니라 어떤 사건이 어느 종목에 언제 중요한지도 판단해야 한다. 이를 위해서는 텍스트, 종목 연결, 거래 세션 시각과 시세를 함께 다루고, 미래 사건이 모델 선택에 들어가지 않는 평가가 필요하다. FNSPID는 미국 종목 뉴스와 시세를 대규모로 연결하지만 \cite{dong-etal-2024-fnspid}, 한국 시장의 공시 식별자, 거래 세션과 회사 별칭을 그대로 표현하지 못한다. KRX Bench와 FINKRX는 한국 금융 지식과 질의응답을 평가하므로 \cite{son-etal-2024-krx,son-etal-2025-finkrx}, 문서 단위 시장반응 분류와는 과제가 다르다.

본 연구의 기여는 다음과 같다. 첫째, 뉴스·공시 550,662건과 전 종목 일별 시세 10.69M행을 비공개 DB가 아닌 검증 가능한 Parquet 파일로 배포한다. 둘째, 발표 세션, 사건 중복, 종목-거래일 충돌과 시간 경계 누수를 통제한다. 셋째, 감성·의미 중요도·사후 시장반응을 서로 다른 과제로 정의한다. 넷째, 강한 기준선, 3개 seed, calibration과 통계 검정, 원천별 실패 분석을 함께 보고한다. 공개 시스템명은 \system{}이다.

\section{관련 연구}

FNSPID는 1999년부터 2023년까지 S\&P 500 기업 4,775개에 대해 뉴스 15.7M건과 시세 29.7M행을 제공한다 \cite{dong-etal-2024-fnspid}. 본 연구는 절대 규모 경쟁보다 한국 공시, 거래 세션 정규화와 사건 혼입 통제를 추가한다. KRX Bench와 FINKRX는 한국 금융 질의응답 및 instruction 평가를 제공하지만 본 연구의 사건-시세 연결 점수와 직접 비교할 수 없다 \cite{son-etal-2024-krx,son-etal-2025-finkrx}. MultiFinBen은 다언어·다중모달 금융 평가를 확장한다 \cite{peng-etal-2026-multifinben}.

금융 감성은 시간 분포 이동에 따라 성능이 저하될 수 있어 \cite{guo-etal-2023-predict}, 무작위 분할보다 시간 외삽 평가가 중요하다. FININ은 뉴스 상호작용과 가격을 함께 모델링하지만 \cite{wang-etal-2024-modeling}, 본 연구는 수익률 최적화가 아닌 해석 가능한 4등급 대리과제를 평가한다. 표현 모델은 한국 금융 말뭉치로 사전학습한 KF-DeBERTa \cite{jeon-etal-2023-kfdeberta}를 LoRA \cite{hu-etal-2022-lora}로 미세조정한다.

\section{데이터셋 구축}

\subsection{원천과 규모}

뉴스는 상장기업 별칭을 기준으로 Naver 검색 결과에서 수집하고, 공시는 OpenDART에서 가져온다. 각 문서는 발표시각, 원천, 제공자, URL 기반 provenance와 정규화 텍스트를 보존한다. 문서-종목 관계는 종목코드와 대표 연결 여부를 포함한다. 시장반응 행은 유효 거래일, 미래 반응 특징, 순서형 라벨, split, 충돌 여부와 사건 군집 ID를 저장한다.

\begin{table}[t]
\centering
\small
\begin{tabular}{lr}
\toprule
구성요소 & 건수 \\
\midrule
뉴스 문서 & 524,696 \\
OpenDART 공시 & 25,966 \\
전체 문서 & 550,662 \\
문서-종목 관계 & 819,772 \\
시장반응 행 & 398,942 \\
비혼입 대표행 & 130,566 \\
일별 시세 행 & 10,691,998 \\
연결된 공시 원문 & 8,972 \\
기본 의미 Gold & 680 \\
\bottomrule
\end{tabular}
\caption{동결 \dataset{} 배포본의 규모.}
\label{tab:scale-ko}
\end{table}

초기 모델 개발에서 사용 가능한 공시 원문은 273건이었다. OpenDART 전체 목록과 원천 식별자를 연결하고 응답 크기 제한, timeout, 재시도, ZIP 검증과 상태코드 기록을 적용해 배포 공시 8,972건의 원문을 확보했다. 전체 공시의 34.55\%이며, 실패 89건은 사유별로 기록했다. 복구 원문은 Gold가 아니다. Gold URL을 제거한 실제 공시 8,302건은 의미 중요도 학습 후보로만 사용하고 감성 정답으로 사용하지 않는다.

\subsection{Gold와 주석}

기본 운영 Gold는 뉴스 80건과 공시 600건이다. 공시 중요도 분포는 LOW 250, MEDIUM 34, HIGH 287, CRITICAL 29이며, 감성은 부정 55, 중립 475, 긍정 70이다. 정책 경계를 검증하는 공시 stress set 310건을 더해 의미 중요도 평가는 총 910건이다. 학습과 Gold의 URL 중복은 0건이다.

코드북은 등급 우선순위, 증거 필드, 감성 override, 모호 사례 제외와 종목별 상한을 고정한다. 상장폐지, 횡령·배임, 감사의견 거절, 부도·회생·파산과 같은 존속위험만 CRITICAL로 제한한다. 단순 거래정지나 불성실공시 지정은 HIGH다. 단, Gold는 공개 코드북을 따른 단일 Codex 기반 검수자가 판정했다. 독립 금융전문가 다중 주석이나 평가자 간 일치도가 없으므로 정책 보강 결과를 외부 SOTA로 해석하지 않는다.

\subsection{거래일과 사건 정규화}

장 마감 전 문서는 같은 유효 거래일, 마감 후·주말·휴일 문서는 다음 거래일에 배정한다. 같은 사건의 반복 보도는 유효 거래일과 정규화 내용을 포함한 지문으로 묶는다. 같은 종목·거래일에 서로 다른 사건 군집이 둘 이상 존재하면 혼입으로 표시해 기본 학습에서 제외한다. 같은 사건·종목·거래일 안에서는 정보량이 가장 큰 문서를 대표로 선택한다.

\section{과제와 라벨}

\paragraph{감성.} 부정·중립·긍정 언어를 예측한다. 균형 공개 시험셋은 등급별 311건, 총 933건이며 운영 뉴스·공시 Gold로 원천 이동을 확인한다.

\paragraph{의미 중요도.} 공시 의미가 기업에 주는 중요도를 LOW, MEDIUM, HIGH, CRITICAL로 분류한다. 시장의 예상 여부와 무관하게 안정적이어야 한다.

\paragraph{사후 시장반응.} 수정주가와 동일 시장 평균수익률의 차이를 1·3·5거래일에 계산한다. 점수는 1일 절대 초과수익률 0.50, 3일 절대 초과수익률 0.20, 60일 로그 거래량 z-score 충격 0.15, 20일 일중 변동성 충격 0.15로 합성한다. 임계값 0.20, 0.45, 0.75로 네 등급을 만든다. 미래 관측은 정답 생성에만 사용한다. 이 값은 인과적 중요도나 거래 신호가 아니다.

\section{누수 통제 실험}

비혼입 대표셋은 Train 107,175건, Validation 6,975건, Test 10,750건이다. Validation은 2026년 1월 1일, Test는 2026년 4월 1일부터 시작하며 각 경계 직전 7일을 embargo한다. 전처리, 임계값, 사전확률 보정, 조기 종료와 모델 선택은 Test를 사용하지 않는다. 통계 비교는 정렬된 문서 ID와 정답이 같은 경우에만 수행한다.

감성 모델은 고정 revision의 \texttt{kakaobank/kf-deberta-base}에 LoRA rank 16, $\alpha=32$, dropout 0.1을 사용한다. 배포 감성은 KF-DeBERTa와 TF-IDF logistic 확률을 0.8:0.2로 결합한다. Validation으로 학습한 stacker도 평가했지만 독립 시험과 운영 Gold가 낮아 배포하지 않았다.

의미 중요도는 제목, 제목+요약, 제목+요약+전문의 TF-IDF logistic을 Validation으로 비교한다. 제목+요약과 제목의 macro-F1은 0.9894로 같지만 Brier score가 각각 0.00134와 0.00156이어서 제목+요약을 선택했다. 전문 모델은 0.9815로 낮아 강제 사용하지 않는다. 존속위험 문구에는 좁은 결정적 floor만 적용하고 모델 단독과 정책 보강 결과를 분리한다.

시장반응 기준선은 문자 2-5 gram TF-IDF와 class-balanced one-vs-rest logistic이다. 신경망은 KF-DeBERTa LoRA rank 16, 256 token, 유효 batch 32, learning rate $2\times10^{-4}$를 사용한다. loss는 class-balanced focal cross entropy($\gamma=1.5$), 가중치 0.30의 ordinal CDF MSE와 label smoothing 0.02를 결합한다. seed 17, 42, 73을 실행하고 Validation macro-F1만으로 배포 seed를 선택한다.

지표는 accuracy, macro-F1, quadratic weighted kappa(QWK), 순서형 MAE, 15-bin ECE와 다중분류 Brier score다. 2,000회 paired bootstrap, exact 양측 McNemar와 70개 시험 거래일을 단위로 한 2,000회 군집 bootstrap을 사용한다.

\section{결과}

\subsection{감성분류}

\begin{table}[t]
\centering
\small
\begin{tabular}{lrr}
\toprule
모델 & Accuracy & Macro-F1 \\
\midrule
TF-IDF logistic & 0.4780 & 0.4423 \\
KR-FinBERT-SC & 0.7363 & 0.7272 \\
KF-DeBERTa LoRA & \textbf{0.8853} & \textbf{0.8850} \\
80:20 ensemble & 0.8842 & 0.8840 \\
Logistic stacker & 0.8821 & 0.8820 \\
\bottomrule
\end{tabular}
\caption{균형 공개 감성 Test 933건.}
\label{tab:sent-ko}
\end{table}

KF-DeBERTa는 KR-FinBERT-SC보다 macro-F1이 0.1579 높다. paired bootstrap 차이는 0.1566, 95\% CI는 [0.1262, 0.1877]이고 exact McNemar $p=1.51\times10^{-18}$이다. 80:20 ensemble은 공개 시험 점수가 근소하게 낮지만 뉴스 Gold에서 accuracy 0.9000/macro-F1 0.8642, 공시 Gold에서 0.9233/0.8344를 기록해 운영 견고성을 위해 유지한다.

\subsection{공시 의미 중요도}

제목+요약 모델 단독은 기본 공시 Gold 600건에서 accuracy 0.9850, macro-F1 0.9470, CRITICAL F1 0.8163이다. 기존 v3 이전 통합 분석기의 0.7150/0.4489/0.2619보다 높다. 좁은 코드북 floor를 더하면 기본셋 1.0000, 기본+stress 910건에서 accuracy 0.9989, macro-F1 0.9962, QWK 0.9994를 기록한다. 이는 같은 코드북으로 만든 floor와 평가셋의 내부 일치도이므로 독립 외부 타당성이 아니다. 모델 단독 결과를 기본 과학 결과로 해석한다.

\subsection{시간 외삽 시장반응}

\begin{table}[t]
\centering
\small
\begin{tabular}{lrrr}
\toprule
모델/실행 & Acc. & M-F1 & QWK \\
\midrule
TF-IDF Train 전용 & 0.4643 & 0.3429 & 0.3141 \\
KF-DeBERTa s17 & \textbf{0.5190} & 0.3724 & \textbf{0.4704} \\
KF-DeBERTa s42 & 0.5030 & \textbf{0.3927} & 0.4627 \\
KF-DeBERTa s73$^\dagger$ & 0.5095 & 0.3820 & 0.4694 \\
3-seed 평균 & 0.5105 & 0.3824 & 0.4675 \\
\bottomrule
\end{tabular}
\caption{시장반응 Test 10,750건. $\dagger$ Validation만으로 선택.}
\label{tab:impact-ko}
\end{table}

선택 seed 73은 기준선보다 macro-F1이 0.0391, QWK가 0.1554 높다. 거래일 군집 bootstrap 95\% CI는 accuracy [0.0351, 0.0557], macro-F1 [0.0256, 0.0536], QWK [0.1331, 0.1776]으로 모두 0을 넘고, McNemar $p=1.70\times10^{-20}$이다. Brier는 0.6450에서 0.6129로 개선되지만 ECE는 0.0491에서 0.0662로 악화되어 calibration이 일관되게 개선되었다고 주장할 수 없다.

\subsection{원천별 실패 분석}

\begin{table}[t]
\centering
\small
\begin{tabular}{llrr}
\toprule
원천 & 모델 & M-F1 & QWK \\
\midrule
\multirow{2}{*}{뉴스 10,160} & TF-IDF & 0.3436 & 0.3220 \\
 & KF-DeBERTa & \textbf{0.3847} & \textbf{0.4779} \\
\multirow{2}{*}{공시 590} & TF-IDF & \textbf{0.3006} & \textbf{0.0611} \\
 & KF-DeBERTa & 0.2211 & 0.0219 \\
\bottomrule
\end{tabular}
\caption{원천별 시장반응 Test.}
\label{tab:source-ko}
\end{table}

전체 개선은 뉴스에서 발생한다. 공시에서는 Transformer가 LOW를 과대예측하며 macro-F1이 0.3006에서 0.2211로 하락한다. 공시 CRITICAL은 10건뿐이고 recall은 0이다. 따라서 공시 가격반응 SOTA를 주장하지 않으며, 공시 전용 장문 모델이나 mixture-of-experts가 후속 과제다.

\section{Ablation과 운영 비교}

TF-IDF에서 원천 prefix를 제거하면 macro-F1은 0.3429에서 0.3427이 된다. 제목+요약만 쓰면 0.3505로 오르므로 약하게 정렬된 전문을 더하는 것만으로 선형 모델이 좋아지지 않는다. 학습량 20k, 50k, 107,175건에서 점수는 각각 0.3260, 0.3446, 0.3429로 단조 증가하지 않는다. 데이터 양만이 병목은 아니다.

v3 전환은 기존 운영셋 일부에서 회귀한다. 실제 뉴스 감성 accuracy는 0.9750에서 0.9000, stock-review 감성은 0.9180에서 0.7880, 실제 뉴스 중요도는 0.9625에서 0.9500, stock-review 중요도는 0.8480에서 0.8060, 사건 macro-F1은 0.7307에서 0.7110으로 하락한다. 과거 라벨 정책과 새 과제가 일부 충돌하지만 이 회귀를 숨기지 않고 과제별 gate와 fallback을 유지한다.

\section{논의와 결론}

\dataset{}의 주된 기여는 예측 과제의 완결보다 한국 금융 데이터 인프라와 평가 규율이다. KF-DeBERTa는 같은 Test에서 감성과 전체 시장반응을 개선하지만 공시 시장반응에는 일반화하지 못한다. 또 의미 중요도와 가격반응을 API에서 분리해야 한다. 파산 공시는 시장이 예상했더라도 의미상 중요하고, 확인되지 않은 소문은 큰 가격변화를 만들더라도 신뢰할 중요 사건은 아닐 수 있다.

FNSPID보다 절대 문서 수는 적지만 한국 공시, 세션 정규화, 사건 혼입 필드, embargo split과 파일 배포를 추가한다. KRX Bench와 FINKRX보다 과제 범위는 좁고 사건 연결에 집중한다. 따라서 교차 benchmark SOTA나 수익성은 주장하지 않는다. 동일 Test의 명시된 기준선 대비 개선, 그리고 공개 코드북에 대한 내부 일치만 방어 가능한 주장이다.

\section*{한계}

Naver 검색 결과는 질의, 제공자, 중복 제거와 생존 편향의 영향을 받는다. OpenDART 전문 커버리지는 34.55\%이며 결측은 무작위가 아니다. 회사 별칭은 사명 변경, 약어와 상장폐지 종목을 놓칠 수 있다. 일봉은 장중 반응, 발표 전 정보 누출과 거시·산업·동시 사건 혼입을 제거하지 못한다. 시장반응은 상관적 약지도다.

기본 Gold 680건과 stress 310건은 전체 말뭉치보다 작다. 특히 공시 Gold는 단일 Codex 기반 판정자가 고정 코드북으로 검수했다. 독립 금융전문가 2인 이상, 합의 과정과 평가자 간 일치도가 없다. 시장반응 Test는 유효 거래일 70개이며 뉴스가 대부분이다. 한국어와 2026년 경계에 한정돼 지역·기간 일반화도 제한된다.

\section*{윤리적 고려}

공개 뉴스·공시·시세도 저작권, 개인정보와 맥락 보존 문제가 사라지는 것은 아니다. 배포는 연구 목적 정책, provenance, 삭제와 버전 절차를 따른다. API key, credential과 비공개 사용자 데이터는 제외한다. 기업 사건에 포함된 개인 식별자는 downstream에서 접근 통제와 목적 제한이 필요하다.

금융 NLP는 소문, 시세조종과 자동화 편향을 증폭할 수 있다. API는 의미 중요도와 관측 시장반응을 분리하고 confidence와 모델 버전을 반환하며, artifact hash나 승격 gate가 실패하면 검증된 fallback으로 fail closed한다. 자동 주문이나 개인화 투자 조언을 수행하지 않으며 고위험 알림은 원문 확인과 사람의 판단이 필요하다.

생성형 코딩 보조 도구는 데이터 엔지니어링, 코드 검토, 코드북 기반 Gold 검수와 원고 작성에 사용됐다. 모든 수치는 버전이 고정된 기계 판독 보고서와 대조한다. 이 사용은 서로 연관된 주석·서술 오류를 만들 수 있으며 독립 전문가 검증을 대체하지 않는다.

\section*{재현성}

환경은 Python 3.12와 lockfile을 사용한다. base model revision, LoRA 설정, seed, optimizer, scheduler, 완료 epoch와 Validation 선택값을 보고서에 저장한다. Parquet은 byte size와 SHA-256을 manifest에 기록한다. 기준선과 Transformer seed별 prediction 파일로 재학습 없이 paired 평가가 가능하다. remote model code는 비활성화하고 safetensors와 artifact version·path·size·hash를 확인한 뒤 로드한다. 논문 표의 핵심 수치는 빌드 전에 동결 JSON과 자동 대조한다.

\appendix
\section{배포 파일과 승격 조건}

배포본은 \texttt{documents.parquet}, \texttt{document\_entities.parquet}, \texttt{market\_impacts.parquet}, \texttt{prices\_daily.parquet}, \texttt{splits.parquet}, \texttt{annotations.parquet}으로 구성한다. 애플리케이션 DB에서 데이터셋을 재구성하지 않는다.

감성 승격은 공개 Test macro-F1 0.85 이상, 동일 ID에서 KR-FinBERT-SC 이상, 뉴스·공시 운영 Gold 각각 accuracy 0.90 이상을 요구한다. 시장반응은 Validation 전용 선택, 기준선과 Test ID 일치, accuracy·macro-F1·QWK의 거래일 군집 bootstrap 하한 0 초과와 artifact 무결성을 요구한다. 전체 gate를 통과해도 원천별 회귀는 차단 진단으로 남긴다.

\begingroup
\renewcommand{\emph}[1]{#1}
\bibliographystyle{acl_natbib}
\bibliography{references}
\endgroup

\end{document}
